\typesetchapter{Chapter 1}{Ascension}

\noindent \textit{A Candidate learns where light falls.}\par
\noindent \textit{A Veil learns who placed it.}\par
\noindent \textit{A Keeper learns what must remain dark.}\par
\noindent \textit{Do not rush to the third sentence because the first wounds your pride.}\par
\noindent \textit{--- Va-Raedin, Veil of Light, spoken during a first morning of instruction}\par

\ornament

He looks to the window before he speaks, every morning, before he will give me a single word. I let him. There is nothing out there but the valley and the cable-car and a great deal of snow, but he studies the glass the way I would study a seismograph --- as though the surface were the visible end of something larger moving underneath.\par

``There it is again,'' he says. ``Your mountain. It is the wrong kind. Made of snow. And it carries people to the top in little glass boxes, and brings them down again, unharmed, in time to eat.'' He watches a pair of the cabins climb the far slope. ``What are they.''\par

``A lift,'' I say. ``A cable, and cars hung from it. It carries people up so they may stand in the snow, then carries them down again. That is all it is.'' There is no percentage in lying to a man about a ski-lift; and I was told to answer what he sees and colour none of it, which is the rare instruction in this business I have not wanted to argue with.\par

``To stand in the snow,'' he repeats, filing it somewhere. ``And be returned.'' He turns from the window to the machine, and then to me. ``Yours announces itself,'' he says. ``Mine spent a thousand years learning not to be found.''\par

I do not answer it. He permits that; we have arrived, over four mornings, at a serviceable working silence.\par

``You speak my tongue well,'' he says --- not for the first time, and still with the same faint disbelief, as though each morning he expects to find the trick of it worn through.\par

``We call it Noxis.'' I keep it flat. ``Four years I was taught it, in a room, by people who would not say where they had kept it. You are the first living soul I have spoken it to.'' It is all true, which is why I say it; I was permitted a short list of true things and told to keep to it, and I have found the list easier to keep than they perhaps expected. ``You are my first task.''\par

``Noxis.'' He turns the word in his mouth like a coin struck in a country he has never walked in. ``You have given it a name. We only ever spoke it.'' Then the fact-checker surfaces under the wonder. ``And your own tongue --- what do you call that?''\par

I give him the name of my language and nothing around it --- not the others, not how many, not that an educated child at home carries two or three of them before he is ten. I am here to take his world down onto the reels, not to hand him mine; he grasps the etiquette without being taught it, and lets a withheld thing lie where I leave it, which is more courtesy than I have been shown by men who share my century.\par

``Paul.'' He says my name the way he sets everything down, exactly, so he can find it later. ``The people I was sent to. The ones the work is for. Are you one of them.''\par

``I am,'' I say. So far as I have been told, it is the entire content of what I am for.\par

He nods, satisfied, and I let the reels turn and keep my face in order. Consider the expense of me. They brought me across the Atlantic to a house I am not allowed to leave, paid me past the whole earnings of a career, and drilled me for four years in a language that on any sober account does not exist --- and the entire outlay, they tell me, resolves to this: the man across the table, and whatever he has come to say. He asks whether I am one of his people as though it were a trifle. I say yes, because it is; and I keep to myself the only part that signifies, which is that I have not yet decided whether a single word of what he is about to tell me is true --- and that an obscene sum has been spent to make me the man who rules on it.\par

``Then begin where it begins,'' I say.\par

\ornament

There is no beginning. You keep asking me for one, Paul, and there is no such thing; the beginning is always one chapter too far back, and I was taught to say so by people who did not want beginnings examined. So I will give you a morning instead. The morning the Mountain opened my hand and told me my name, and I was young enough to believe that being named was the same as being freed.\par

You have a machine on the table between us. A small wheel of tape that turns when I speak and stops when I stop, and you have told me, twice now, that it is only so you will not have to trust your memory. I believe you. But understand what you are asking of me. I come from a place where memory was an office, guarded and rationed and punished when it wandered. And here you sit, letting a recorded voice run as though sound were a safe thing. As though a voice could not be carried out of a room and used against the throat that made it.\par

That is the first thing you should learn about where I come from, before any of the rest. Not the darkness. Not the stone. This: that we did not fear seeing. We feared \textit{carrying}. A man could look upon anything, so long as he left it where it lay. The crime was never the eye. The crime was the hand that took the thing away.\par

I was fifteen. Hold on to that. Everything I am about to describe, I saw with the eyes of a boy who believed that brightness was a kind of virtue.\par

\ornament

The Hall of Ceremony had been prepared for three days. That was the word they used, \textit{prepared}, for the whole patient business of lowering lamps and filling bowls and polishing dull stone into a greasy obedience. But the Hall had not truly been prepared in three days. It had merely been cleaned for us. The real preparation had taken centuries and left its illness in the rock: an enormous upward hollow mined into the Mountain, round at the floor, narrowing as it climbed, never quite a cone, never quite anything a person could name without lying a little. And high above, past the pillars and their dark ring of stone, the throat from which the Thread came down.\par

The Thread was thin as if it had no right to hold what it held, and at the end of it hung the Pendulum.\par

Forget whatever picture you are making. It was not a ball. It was a thick round plate, black-rimmed, wider than a man is tall, two-faced, solemn in the idiotic way that very heavy things are solemn, because they need no anger to be dangerous. One face bore God, hard-cut and dry-eyed, less a comfort than a verdict. The other was the Beast, though we were not encouraged to say so with appetite, because children love terrible names and adults dislike hearing their own fears enjoyed. As the disk crossed the Hall it turned, showing God to one side, then edge, then the Beast to another --- a slow exchange of blessings and threats performed by metal and by weight.\par

At the top of the Thread the lean was almost nothing. That was the obscenity of the scale. A small indecency away from the vertical, up there; but below, at the end of a hundred and twenty-five metres of drop, that little lean became a true crossing, five metres and more from the centre before the disk came back through the middle and went to the other side. And the road it travelled did not stay put. Swing after swing it seemed the same patient menace, back and forth, back and forth --- but the line of its travel turned, slowly, across the floor, until after nearly three days it had made a full round and come back to itself, as though the Hall had spent all that time chewing one invisible circle.\par

\ornament

``I know the object you are describing,'' I say, when he has done with the disk. He stops, mildly put out that I should. ``We have them. A Foucault pendulum --- a heavy bob on a long wire, swung once and left to it. The plane of the swing appears to rotate across the floor over the better part of a day; in fact the floor is turning, which is to say the Earth is, and the pendulum only keeps to the fixed stars.'' It is the single thing he has said all morning that I can put my hands on, and I take an unprofessional comfort in it. ``I promised you at the outset I would say when I did not believe a thing,'' I go on, ``and I do not believe most of what you have told me this morning. But the pendulum you have just described is real, and you have described it correctly, which is a considerable inconvenience to me. We hang the largest of them in halls of glass, behind a rope so no one interferes, with a small printed card to spare the visitor the labour of understanding. We send schoolchildren to watch. They are bored, and afterwards examined on it.''\par

He is very still. ``A hall of glass,'' he says. ``With a rope. So that it is seen and not touched.'' Something moves behind his eyes that I do not name for him. ``What is the hall called.''\par

``A museum,'' I say. ``A building where a people store what they have finished using, so the young may file past and be assured they came from somewhere.''\par

``A house for things you have finished with,'' he says slowly. ``And you rope them off, and explain them on a card, and send your children to be bored by them.'' He does not finish the thought, and I decline to finish it for him: that his people knelt to that swinging weight in the dark and assigned it to God, and mine put the identical weight behind a rope and assigned it to the rotation of the Earth --- the same pendulum, in both halls. He holds the look a moment longer than I care for. I move him on before it can become a question --- not because I am hiding an answer, but because I have not been sent here to supply one.\par

``Go on,'' I say. ``You had knelt.''\par

\ornament

Around the cleared and forbidden centre the rest of us had been packed into the Hall's permitted flesh. Five pillars stood in a ring, one for each of the charges by which our whole world was divided --- Balance, Faith, Shield, Light, Justice --- each so crowded with carving that the stone seemed to have swollen with symbols like a diseased thing. I stood beneath Light, my side, opposite the shut Mountain gates and the black mass of the old gravure, and I held my shoulders as though the entire Hall had assembled to judge their angle. The robe scratched my throat. The lamps made the air taste of oil. My feet had begun their own small, exact, deeply untheological suffering. The ceremony had not even begun.\par

Beside me a boy kept wetting his lips with scripture, not from belief, but the way frightened fingers worry a scab.\par

\begin{scripture}
Beyond the forest thou shalt not wander.
\end{scripture}

He had said it three times. I did not turn toward him. Turning in the Hall was never simply turning; it made cloth move and sandals complain and some adult eye discover that your body had not yet been entirely surrendered. So I faced forward, toward the disk, toward everything I did not understand and was expected to honour in the meantime.\par

Then a single chime entered the Hall already enlarged by stone, one clear note scraped clean of mercy, and the almost-cone answered it in climbing layers. Heads lowered. Robes settled. And the Faith chorus began to sing --- a long hymn, and I knew it would be a long and tiring ceremony --- and near the end of it came the line that was the whole law of my people folded into six words.\par

\begin{scripture}
Let every footfall be measured,\\
Let every breath be whole;\\
Let no one hand hold the whole.
\end{scripture}

\ornament

``Let no one hand hold the whole,'' I say.\par

He looks up from the machine. I had not meant to say it aloud; it simply arrested me, the way a clean statement of a principle will. I let the lapse pass without comment and rather hope he will do the same.\par

``You know the line,'' he says. It is not quite a question.\par

``It was in the papers they gave me to study,'' I say. Which is true, and as much as I intend to give him.\par

He watches me for a moment, the way he watches the window. ``You will want me to explain it,'' he says. ``I can see it in your face. But if I explain it now, at the start, I spoil the only thing worth telling you --- which is how long it took me to understand what it cost. Let it stand for now as a song. It will become a wound later. Everything in the Mountain did.''\par

I do not press him. If there is a method to what I do here, it is that a pressed account is a contaminated one --- I would sooner have his order of things than impose my own. So I nod, and mind the reels, and let the silence stand in for agreement. He has decided the line is a wound waiting to open. I have decided nothing yet; declining to decide prematurely is, in fact, the precise thing they are paying me for.\par

\ornament

What we had gathered to witness, before my own small part, was a death that was not permitted to be a death.\par

The Keeper of Faith was being replaced. Ka-Mereth --- old past telling, eaten from within by age or prayer or office, her white cowl slipping from a skull too small for it --- was carried in on a low chair of dark wood, and the Hall made the small animal movement a crowd makes when it has promised to be still and is suddenly given something dying to look at. She was not dead. That was the trouble. Death would have been cleaner.\par

Behind her walked the woman who would take her charge, in a pilgrim's mantle mended in darker patches, dragged through stone and sweat until it no longer quite belonged to cloth. At her throat two raw lines crossed the skin. Not ornaments. Not symbols yet. Wounds first, red and badly healed, which meant the Hall could make them sacred only after every one of us had been made to see that they had hurt. I stared at them. I was fifteen; I stared at wounds, at women, at authority, at anything the adults had not yet taught me how to look at without giving myself away.\par

And here Light gave me my first lesson, though I did not recognise it as one. Va-Raedin, who stood behind me, touched two fingers to the back of my wrist. He did not even look at me. That was the cruelty of it --- no rebuke, no lifted eyebrow, nothing splendid to resent, only two fingers laid against my wrist with the hygienic delicacy of a man removing a stain from glass. I lowered my eyes, and the whole Hall shrank for a moment to the disgrace of having been caught feeding on another person's pain as though the ceremony had served it to me.\par

Then came the thing I most want you to hear, Paul, because it is the thing your world would refuse to believe, and it is true.\par

They brought out a bowl.\par

The Bowl of the Last Prayer, they called it, and it had not been seen in a lifetime, so that a murmur escaped the crowd before discipline could throttle it. It sat in a black cradle, pale and elliptical, its outside tortured into sanctity --- five bands of scripture, hooded figures no larger than fingernails, gates and roots and chains and waves cut into the alabaster until no innocent surface remained. But the inside had been spared. After all that outer delirium the interior opened into a long clean hollow, polished to a milky silence, bare of any word. Faith had decorated everything except the thing that mattered. The outside belonged to worship. The inside belonged to water. And rising from the pale floor at either end stood two little columns of old black bronze, no taller than a finger-joint. The Ear, and the Mouth.\par

The dying Keeper was carried close enough to look inside. Her hand shook so violently that two attendants moved to help, and she flicked them away with a dry remnant of a gesture, and they obeyed. She lowered a narrow rod toward the Ear. Then her hand dropped the last finger's width, and the sound was almost nothing --- a small dry tick of metal on metal, miserable after so much expectation. I remember the shame of having expected thunder.\par

Then the water moved.\par

The first rings left the Ear, widened, touched the clean curve of the alabaster, and came back \textit{wrongly}. They ran along the pale wall and bent inward and crossed themselves, and for one instant they made, near the Mouth, a shape I could not name then and will not pretend to name now. It was not a picture. It was not writing. It was a little treason of water against the common law of the eye. And then the Mouth answered --- a thin, high note, so narrow it seemed less heard than discovered inside the teeth --- and died, and the water smoothed itself with the hateful innocence of a thing that has proved a point and refuses to explain it.\par

The old Keeper's eyes, almost extinguished a moment before, sharpened with a private violence.\par

``The chamber answers,'' she said.\par

I did not understand it. I want to be honest with you about that. I was a boy, and I thought I had watched a miracle. It would be years before I understood that I had watched no miracle at all, but a mechanism; that Faith did not pray so much as it \textit{listened}, that a prayer could be a shaped sound sent through stone to be answered, altered, somewhere else. That everything holy in the Mountain was a machine wearing the robe of God. But that morning I only knew that something below the bowl, something behind the engraved clutter of our ignorance, had replied when it was called. And that no one was going to tell me how.\par

They named the new Keeper then. The great cowl, yellowed and heavy with old thread, was lowered over the kneeling woman like a burden rescued from a corpse, and I saw --- without yet having the words for it --- that office does not merely crown a body. It edits one.\par

``Rise now as Ka-Elun, Keeper of Faith.''\par

She rose. Badly, slowly, her knees stiff, dust falling from one fold of her robe in a little grey sigh --- not at all like the clean ascent the songs pretended. And the old Keeper was carried away alive through the Faith arch, into chambers I would never see, and the children were told such retirement was holy. I thought it looked like being sealed slowly while still breathing, and I hated myself for the thought, which did not stop the thought, and which is how I first learned that shame is not repentance. It is only the body's alarm at having produced knowledge before it had permission.\par

\ornament

Then the ceremony changed, and did not end, and that was the first warning.\par

The great Ascension was complete; the Kahir should have withdrawn. Instead the bowls stayed on their stones, Shield did not release the exits, and behind my shoulder Va-Raedin's hand withdrew --- not in disapproval, but in attention.\par

``Now,'' he said.\par

My mouth went dry.\par

I had known it was coming. My mother had washed my robe herself that morning, despite the attendants Light had sent. My father had held my wrist and told me to stand straight when spoken to. But knowing is a poor instrument when the thing known begins to move toward your body. The naming of the young was not the heart of the day, and that made it worse, because there is a particular humiliation in being told you matter only as evidence that something larger than you has not yet failed.\par

I descended into the Hall, and the Hall became larger with every step. I passed below the Shield arch, and she was there.\par

I will tell you about Va-Sheva properly another time, Paul, or perhaps I never will, because some debts are not improved by being described. That morning it is enough to say she was nineteen and already finished in a way I was not, dark leather over the sleeveless robe, a short blade at the hip and a longer one across her back, standing with the awful efficiency of someone shaped before the rest of us had begun to set. She had once let me follow her through the service tunnels. She had once laughed at me for crying over cut lamp-glass, and then stopped laughing, which was worse. When her eyes found mine for an instant it ruined my balance more surely than any uneven stone, because there are glances that do not give hope but make hope behave indecently --- standing up inside you at the exact moment you need it to kneel.\par

Then I knelt on the eastern stone, and it was warm under my knees, too warm, someone having opened a vent beneath the stone more than the ceremony required, and I remember it because fear attaches itself to small things --- to heat, to the ridge of a sandal, to anything smaller than the question waiting above.\par

Ka-Syphiron, Keeper of Light, watched me with pale eyes that seemed to look not at a thing but through the excuses the thing had prepared in case it was seen.\par

``Do you seek Light?''\par

``I seek Light.''\par

``Why?''\par

That question was not in the rehearsal.\par

I felt the Hall tighten --- not in sound, but in that gathering of attention a crowded body knows before the mind admits it --- and I understood that the safe answers had all been quietly removed from the table before I knew we had left the script. \textit{Because I am worthy} was too proud. \textit{Because I want to know} too naked. And the words were already moving in me faster than caution could inspect them.\par

``Because what is hidden still has shape,'' I said. ``And if it has shape, it may be found.''\par

The Hall went silent.\par

Not the silence of reverence. Not even of shock. A worse silence --- an administrative silence, a silence already making copies of itself for later rooms, because everyone present understood that something had been said which would not remain inside the ceremony. Something small enough to deny and large enough to record.\par

Ka-Syphiron's expression did not change, which frightened me more than anger would have. ``A Candidate does not find all that has shape,'' he said. ``What does a Candidate learn?''\par

The script had returned, and relief made me nearly stupid.\par

``Where light falls.''\par

``And what does a Veil learn?''\par

``Who placed it.''\par

``And what does a Keeper learn?''\par

Here I was to say the thing every boy could recite. But after the torn mantle, after the old woman carried away alive, after the Chamber of the Last Prayer had crossed the Hall as a real place and not a story, the answer no longer sounded like wisdom. It sounded like a door closing.\par

``What must remain dark,'' I said.\par

He lifted his left hand, and Va-Raedin set a blade into it. It was not ornate. That mattered. Light did not decorate the thing that cut. The handle was bone, worn smooth by fingers older than mine, and the edge was bright in a way the lamps did not explain but merely revealed.\par

I extended my hand. There are moments the body understands before the mind, and mine understood that it was being admitted into doctrine not as a witness but as \textit{material} --- the way oil and water and wool and stone were admitted, things to be used because they could be measured and marked. The cut came across the base of my thumb. Quickly. Too quickly for courage. And the blood gathered under the lamplight, darker and thicker than I expected, stupidly private, embarrassingly mine --- though the whole point of the rite was to make it stop belonging only to me.\par

He pressed my thumb to a plate of polished stone marked with five faint circles, so pale I had not seen them until my blood moved toward one and touched only the eastern mark. Only Light. Only the office that had praised my hunger by cutting the part of me that reached. The stone was cold. That is the detail memory kept --- not the words first, but the cold. My blood was warm and the stone was cold, and between them the Mountain made a contract it did not ask me to read.\par

``Blood enters service,'' Ka-Syphiron said. ``Sight enters discipline. Pride enters correction.''\par

I repeated it, because a mouth busy with formula has less room for panic, and because the Mountain knew, better than any physician, that a wounded body will obey a sentence if the sentence is given quickly enough.\par

``Rise now as Sa-Kailan, Candidate of Light.''\par

The name did not fall on me like praise, not warm, not radiant, not like the Mountain opening its arms to what I had always secretly been --- but like a door closing with me on one side of it, and I could not yet tell which side. The Hall did not cheer. Cheering belonged to the stories we told about Outsiders and their kings, noisy people who needed sound to convince themselves an event had happened. The Mountain \textit{acknowledged}. The people struck two fingers against stone, once, twice, five times, and the sound filled the Hall like rain made hard --- not applause, never applause, but agreement turned mineral, a thousand small impacts accepting that the old order had made room for one more wound.\par

I had thought I would feel transformed. Instead I felt cut, watched, and unfinished.\par

Va-Raedin leaned close enough that only I could hear.\par

``You answered poorly,'' he said.\par

My stomach dropped so cleanly that for a breath I forgot the blood. Then he added, ``But not uninterestingly.''\par

That was worse. It was also not entirely bad. Light, I would learn, had a genius for exactly that cruelty --- the wound that also feeds, the insult that opens a door, the correction you store like contraband because it is the nearest thing to affection the office permits. I did not yet know what to do with a thing that could diminish and enlarge me in the same instant. I have spent much of my life since among such things.\par

\ornament

I looked for her again as the Hall began to empty --- Va-Sheva, at the northern arch --- and she was no longer watching the ceremony. Two wardens had come close to her, one speaking without moving his head, the way Shield speaks when the words themselves must not become visible. Her expression tightened. Not fear. Readiness. A hand near the hilt, a glance past the galleries, a small adjustment of stance that changed the space around her. There was almost nothing to see, which is why I remember it so clearly. Shield did not need drama to become Shield. It became Shield by withdrawing mercy from uncertainty.\par

Something had happened, or was about to. I did not know what. I would not know for a long time. But I have learned to trust the memory of that moment more than most, because it was the first time I understood that the ceremony I thought I was the centre of had been, all along, a smaller thing happening in front of a larger one.\par

Then she was gone into the Shield passage, and the arch swallowed her without ceremony, which was proper, because Shield decorated departures no more than it decorated threats.\par

Va-Raedin bound my thumb with a strip of pale cloth. His fingers were efficient. Not gentle, not ungentle --- Light had no use for either when tightness would do. ``You will report to the eastern practice hall before the First hour,'' he said. ``Do not eat heavily. Do not arrive proud.''\par

``I am not proud,'' I said.\par

He tied the cloth too tightly. ``You are fifteen. Pride is your natural condition.'' And when I said nothing, because every path of objection led toward proving him right, he saw that too. ``Good,'' he said. ``You have learned one thing already.''\par

There were men in the Mountain who could praise you and make you smaller. Va-Raedin could correct you and leave you standing taller than you had any right to. I did not understand yet that this is the most dangerous kind of teacher a person can be given.\par

The Hall emptied around me, the lamps dimming one by one as the attendants pinched the wicks, and on the northern wall the old gravure darkened until the carved Mountain lost its edges and the hidden Oasis disappeared into shadow --- so that for a moment the stone showed neither refuge nor promise nor history, only a blackened place where meaning had been. I looked at that shadow until my eyes began to invent shapes inside it.\par

Then I turned toward the eastern passage and followed Va-Raedin into Light.\par

\ornament

I stop the tape. I do not decide to; my hand does it, and the reels coast down, and he marks it, the way he marks everything.\par

``You stopped it,'' he says. He is not troubled. He is almost gentle. ``You do that when you are afraid of what you want to ask next.''\par

I do not tell him what I want to ask. I lean in instead, which he reads correctly as permission, and readies himself to give me the rest of it --- the mechanism under the ceremony, whatever the bowl was calling to, the reason five hands were never allowed to close on one thing at once.\par

``I will tell you all of it,'' he says. ``That is why I was sent. You are Horizon; the account is yours by right, and a thing given to its own people is not carried anywhere it should not go.'' He says it plainly, without weight, the way a man says a thing he has stopped needing to check. ``What my whole world was built to keep from wandering, I will lay in your hands, and you will hold the whole of it, because you are the one place it is safe.''\par

He believes it --- that is the remarkable part, and the part I cannot yet account for. He hands me the one thing his entire world is built to keep from crossing a threshold, and he does it without a flicker of fear, because he is certain he is handing it home.\par

``Start the tape again, Paul,'' he says. ``I have only reached the naming. There is so much more you will want.''\par

I start the tape. The reels turn. My brief, when they set it in front of me, reduced to a single ungracious sentence: prove the man does not exist. So I hunt the seam --- the join where an invention meets itself and fails to line up. Four days, and I have not found one. A poorly made story betrays itself early --- the maker forgets on Friday which lie he told on Tuesday. This one has not slipped once, and the man telling it shows none of the strain a fabrication of this magnitude ought to cost him. I record it as an observation and not a verdict --- \textit{non liquet}, the old courts wrote, when the evidence would not yet resolve: it is not clear. Four days in, \textit{non liquet} is the whole of my professional opinion, and an obscene sum is being paid for it. The account is either true or the most disciplined invention I have ever been handed --- and my part is not to guess which, but to keep the reels turning until the thing settles the question itself.\par

\ornament
